\section{Đánh giá thiết kế (Design Evaluation)}

\subsection{Evaluation Objectives}
% Mục tiêu kiểm thử:
% - Kiểm tra tính dễ dùng và hiệu quả của luồng đặt lịch mới
% - Đo lường mức độ rõ ràng của tính năng theo dõi tiến độ và dặn dò dị ứng

\subsection{Evaluation Method}
% Phương pháp kiểm thử:
% - Usability Testing kết hợp giao thức Think-aloud
% - Bảng câu hỏi đo lường mức độ hài lòng (SUS - System Usability Scale)

\subsection{Participants}
% Số lượng và tiêu chí tuyển chọn người tham gia kiểm thử (ví dụ: 5 chủ nuôi thú cưng)

\subsection{Tasks}
% Danh sách các task kiểm thử:
% - Task 1: Đặt lịch tắm & cắt tỉa cho thú cưng có tiền sử dị ứng
% - Task 2: Kiểm tra và theo dõi tiến độ chăm sóc trên ứng dụng
% - Task 3: Xem lại lịch sử dịch vụ và ghi chú của lần chăm sóc trước

\subsection{Metrics}
% Các chỉ số đo lường:
% - Task Completion Rate (%)
% - Time on Task (giây)
% - Error Rate / Number of Errors
% - SUS Score

\subsection{Procedure}
% Mô tả quy trình tổ chức buổi kiểm thử

\subsection{Results}
% Trình bày kết quả đo lường thực tế qua bảng biểu và biểu đồ

\subsection{Discussion}
% Thảo luận ý nghĩa kết quả:
% - Những điểm cải thiện vượt trội so với quy trình cũ
% - Những vấn đề còn tồn đọng người dùng gặp phải trong quá trình test
